\section{Limits of the method}
\label{sec:limits}

The saving \(\SstarLo\) is the end of what the parameters give.  This section
records why the base identity cannot be exchanged for a better one of the same
kind, and what would have to change.  The comparison of the shapes of
Sch\"onhage's family is certified (Proposition~\ref{prop:shapes}): it is a
statement about the saving formula of Sections~\ref{sec:exact}
to~\ref{sec:optimum} evaluated for other shapes, since the programs exist for
the shape of the paper only.  The rest of the section, the mixtures, the
hypothetical shapes, the experiments and the remark on \(\omega\), is \nf; its
numbers are computed in floating point by the scripts in the directory
\texttt{search} of the repository.

\subsection{The family of Sch\"onhage}

For \(k,n\ge2\), Sch\"onhage's identity for
\(\schon k1n\oplus\schon1{(k-1)(n-1)}1\) has \(kn+1\) terms \cite{Sch81}.  It
has \(s=kn\) outer outputs with one term each, and one shared term for an
inner product of length \(b=(k-1)(n-1)\).  The paper uses \(k=n=3\): \(s=9\),
\(b=4\), ten terms.  The analysis of Sections~\ref{sec:exact}
to~\ref{sec:optimum} is written for this shape.  For another member of the
family it would be repeated with the corresponding counts of terms, variables
and leaves, an asymmetric shape being used together with its transpose so that
tiles stay square; we do not write the programs for other shapes, and what
follows is the saving formula of Theorem~\ref{thm:sup} evaluated for them.
With \(s\) outer outputs and inner length \(b\) one has \(D=b^{m}\),
\(N_0=\sqrt s^{\,L-m}\),
\(\gtwo\) with \(\delta\ln s\) in place of \(\delta\ln 9\), \(\gX\) and \(q\)
with \(\ln b\) in place of \(\ln 4\), the window \((0,s/(s+1))\) for \(\theta\)
in place of \((0,0.9)\), and the thinness constant
\(\Apr(c)=\frac c2\Hent(\frac1c)+(c-1)\frac{\ln s}2\); the saving at the ratio
\(c>s+1\) is then \(S(c)=F_{\Apr(c)}(\Gam(c))\) as in Theorem~\ref{thm:sup}.
In Lean (directory \lean{Optimum/Shapes}): \lean{g2S}, \lean{gammaXS},
\lean{qS}, \lean{gammaQS}, \lean{GamS}, \lean{thinRS}, \lean{savingS},
\lean{basePrunedS} and the saving \lean{shapeSaving}~\(s\,b\,c\); for
\((s,b)=(9,4)\) they reduce to the functions of Sections~\ref{sec:exact}
to~\ref{sec:optimum} (\lean{shapeSaving_nine_four}, from \lean{g2S_nine},
\lean{gammaXS_nine_four}, \lean{qS_nine_four}, \lean{GamS_nine_four} and
\lean{basePrunedS_nine}).  Lemma~\ref{lem:crossing} and the monotonicity of
\(\Apr\) carry over under \(2\le b<(s+1)^{2}\) and \(c>s+1\)
(\lean{GamS_pos}, \lean{GamS_lt_two_thirds}, \lean{GamS_monotoneOn},
\lean{basePrunedS_monotoneOn}).  Table~\ref{tab:family} gives the optimum for
each shape, with pruned encodings and the prime at its fixed point.

\begin{table}[ht]
\centering\small
\setlength{\tabcolsep}{4.5pt}
\begin{tabular}{lrrlrrr}
\toprule
shape \((k,n)\) & \(s\) & \(b\) & \(\sup_{c}S(c)\) (certified)
  & best \(c\) & \(\Gam\) & \(\varepsilon\)\\
\midrule
\((3,3)\) & 9 & 4 & \((\SstarLo,\SstarHi]\) & 21.92 & 0.0759 & 0.0552\\
\((2,4)\) & 8 & 3 & \((0.001875,0.001877]\) & 20.11 & 0.0749 & 0.0501\\
\((3,4)\) & 12 & 6 & \((0.001867,0.001869]\) & 27.96 & 0.0746 & 0.0501\\
\((2,5)\) & 10 & 4 & \((0.001776,0.001778]\) & 24.24 & 0.0742 & 0.0479\\
\((2,3)\) & 6 & 2 & \((0.001669,0.001671]\) & 16.03 & 0.0742 & 0.0450\\
\((2,6)\) & 12 & 5 & \((0.001623,0.001625]\) & 28.39 & 0.0733 & 0.0443\\
\((3,5)\) & 15 & 8 & \((0.001613,0.001615]\) & 34.04 & 0.0732 & 0.0441\\
\((4,4)\) & 16 & 9 & \((0.001573,0.001575]\) & 35.96 & 0.0729 & 0.0432\\
\((4,5)\) & 20 & 12 & \((0.001323,0.001325]\) & 44.01 & 0.0714 & 0.0371\\
\((5,5)\) & 25 & 16 & \((0.001098,0.0011]\) & 54.02 & 0.0698 & 0.0315\\
\bottomrule
\end{tabular}
\caption{The saving for Exact Triangle with Sch\"onhage's identity of shape
\((k,n)\) and pruned encodings.  The enclosures of the suprema over the ratio
\(c>s+1\) are certified in Lean (Proposition~\ref{prop:shapes};
Theorem~\ref{thm:global} for \((3,3)\)).  The last three columns, the ratio at
which the supremum is approached and the parameters there, are computed in
floating point and are not certified.}
\label{tab:family}
\end{table}

\begin{proposition}[The family of Sch\"onhage]\label{prop:shapes}
For each of the nine shapes of Table~\ref{tab:family} other than \((3,3)\),
the supremum over \(c>s+1\) of the saving \(S(c)\) of the shape lies in the
interval given in the table; in particular it is below \(0.00188\).  For
\((3,3)\) the saving exceeds \(\SstarLo\) at \(c=\cCond\).
\end{proposition}

\begin{proof}
As for Theorem~\ref{thm:global}, shape by shape: a lower bound is the value at
one rational ratio, and the upper bound is certified on cells covering
\((s+1,c_{\max}]\), with \(c_{\max}\) between \(200\) and \(600\), and beyond
\(c_{\max}\) by \(\Gam(c)<2/3\) and the monotonicity of \(\Apr\); the cells are
finest on the window around the maximum, which is \([19.6,20.6]\) for
\((2,4)\) and \([52.1,56.1]\) for \((5,5)\).
\end{proof}

In Lean: \lean{shapes_table} collects the nine upper bounds and the value for
\((3,3)\); the enclosures are \lean{sSup_shapeSaving_S24_mem} and its
analogues for the other shapes, from the per-shape bounds
\lean{shapeSaving_S24_lt} and \lean{shapeSaving_S24_gt}.  The
\LeanShapesCells{} cells, in \LeanShapesCellFiles{} generated files
\lean{Optimum/Shapes/Cells_S24_1.lean} and so on, are produced by
\texttt{search/certs.py} with the option \texttt{--encoding shapes} and
checked by Lean's kernel as in Theorem~\ref{thm:global}.

The shape \((3,3)\) is the best.  The exponent \(\Gam\) hardly varies over the
family, and \((3,3)\) also has the largest \(\varepsilon\).  For the same
reason mixtures do not help (\nf): when two shapes are used on different
levels, the optimizer drives the share of the second shape to zero, because no
trade-off between \(\Gam\) and \(\varepsilon\) is available.

\subsection{Hypothetical shapes}

This subsection is \nf; its numbers are floating point.  One can ask what an identity with \(s\) outer outputs, one term each, and a
shared term of inner length \(b\) would give, whether or not it exists.  The
saving grows with \(b\) and falls with \(s\).  For \(s=9\) it is \(0.00210\)
at \(b=4\), \(0.00254\) at \(b=5\) and \(0.00293\) at \(b=6\); for \(s=8\) it is
\(0.00188\) at \(b=3\) and \(0.00252\) at \(b=4\).  In the computed grid (\(5\le s\le12\), \(3\le b\le8\)) every cell with a
larger saving than \((s,b)=(9,4)\) is outside the class of Sch\"onhage's
construction: an identity of this kind with \(s=kn\) satisfies the pairing
bound \(b\le(k-1)(n-1)\), so \(s=9\) allows \(b\le4\), \(s=8=2\cdot4\) allows
\(b\le3\), \(s=6\) allows \(b\le2\), \(s=10\) allows \(b\le4\), \(s=12\) allows
\(b\le6\), and \(s=5,7\) have no factorization.  So within the class of
identities with one term per outer output, the identity of the paper is optimal
for this method.  Even the realistic cells beyond the bound
(\(s\in\{8,9\}\), \(b\in\{4,5\}\)) give savings of only \(0.0025\) to
\(0.0031\), that is, 3SUM in about \(n^{1.9985}\).

\subsection{Experiments on the arrangement of the wanted set}

This subsection is \nf.  The bound \eqref{eq:lemma11} is a worst case over the wanted sets.  A random
wanted set of the same size meets it up to lower-order factors: a leaf of
order \(d\) is shared by at most \(\binomial{L-m+d}{d}\) output strings, and the
expected number of visited leaves of order \(d\) is about
\(\min(|U|\alpha_d,\beta_d)\).  A gain needs a wanted set that is close to a
union of sub-tiles.  Simulation cannot test this: feasible sizes have
\(c=L/m\approx3\), where the \(\beta_d\) grow with \(d\) and every wanted set
saturates the tree, while the relevant regime \(c\approx21\) has at least
\(10^{21}\) leaves already for \(m=1\).  Zeros in the encodings do not help
either: an encoding at a leaf with \(\Pz\) on a set \(Z\) of levels is a signed
sum of at least \(3^{|Z|}\) entries, and for the dominant leaves
(\(|Z|\approx0.89m\)) the expected number of nonzero terms tends to infinity.

\subsection{Open problems}

\begin{enumerate}
\item \emph{New sharing structures.}  An identity in which outer outputs are
  fed by several terms, or with several inner outputs, is outside the leaf
  counting of Section~\ref{sec:recursion}.  Scoring such an identity needs a
  generalization of \(\alpha_d\) and \(\beta_d\).
\item \emph{Structured wanted sets.}  The reductions from 3SUM produce Exact
  Triangle instances whose weights are array entries at additively combined
  indices.  Whether the wanted sets are then far from random is not known.
\end{enumerate}

\begin{remark}[\nf]
The exponent of matrix multiplication plays no role.  The choice of the prime
costs \(n^{\omega}D'^{3/2}\), which is of lower order with Strassen's
\(\omega\le\log_2 7\) already; rectangular fast matrix multiplication would not
change any exponent of this paper.
\end{remark}
